\section{Exact classical Euler tails and the remaining simplicity problem}
\label{app:root:scope}

This appendix records the related classical Euler tail program at its
stationary-separation checkpoint of 5 October 2026. Its purpose is to state
the strongest proved reduction and its exact remaining obligation. The
universal simplicity problem below is a separate question from
Nishibuchi's Conjecture~3.9, which the congruence results of this chapter
resolve. It is also separate from the presentation and gluing questions
called Native~13.2 and Native~13.3 in the original Paper~IX.

The tail results have mathematical proofs, supplemented in the indicated
families by complete coefficient certificates and proved bounds for every
omitted coefficient. Bounded sampled values support implementation checks.
They do not supply the universal quantifiers. The latest continuation did
not run Lean. Its mathematical claims must therefore be read with the
formal-coverage distinctions in the Lean guide.

\subsection{The ordinary Euler packet identity}
\label{app:root:identity}

To distinguish ordinary signed Euler numbers from the generalized numbers
of the main text, write
\[
 \operatorname{sech}t=\sum_{n\ge0}\mathcal E_n\frac{t^n}{n!},
 \qquad \frac{t}{e^t-1}=\sum_{n\ge0}B_n\frac{t^n}{n!},
 \qquad B_1=-\frac12.
\]
Throughout this appendix, $v_2(2)=1$ and $v_2(0)=\infty$.
Let $s_2(n)$ be the number of ones in the binary expansion of the
nonnegative integer $n$. Define
\[
 A(N)=\frac{3^{2N+2}-1}{4(2N+1)(2N+2)},\qquad
 \left(\frac z{\sinh z}\right)^2=\sum_{m\ge0}\beta_m z^{2m},
 \qquad \kappa_m=2^{2m}(2m)!\beta_m.
\]
The initial values are $\beta_0=\kappa_0=1$. For $m\ge1$, differentiation
of $z\coth z$ gives
\begin{equation}
 \beta_m=-\frac{2^{2m}(2m-1)B_{2m}}{(2m)!},\qquad
 \kappa_m=-2^{4m}(2m-1)B_{2m}.
 \label{eq:root:beta}
\end{equation}
The classical dyadic Bernoulli denominator law
$v_2(B_{2m})=-1$ and Legendre's formula yield
\begin{equation}
 v_2(\beta_m)=s_2(m)-1,\qquad
 v_2(\kappa_m)=4m-1.\label{eq:root:weights}
\end{equation}
The source identity
\[
 \sum_{N\ge0}A(N)\frac{t^{2N}}{(2N)!}
 =\frac{\sinh(2t)\sinh t}{2t^2}
\]
implies, after multiplication by $4t^2/\sinh^2(2t)$,
\begin{equation}
 \mathcal E_{2N}=\sum_{m=0}^N
    \kappa_m\binom{2N}{2m}A(N-m).
 \label{eq:root:packet}
\end{equation}
For a retained depth $M\ge0$, put
\[
 \mathcal T_M(N)=\mathcal E_{2N}
   -\sum_{m=0}^{\min(M,N)}\kappa_m\binom{2N}{2m}A(N-m).
\]
An empty tail, $N\le M$, is exactly zero.

\begin{proposition}[Sharp uniform classical truncation bound]
For every $M,N\ge0$,
\[
 v_2\bigl(\mathcal T_M(N)\bigr)\ge4M+3.
\]
For each fixed $M$, equality holds at $N=M+1$.
\end{proposition}
\begin{proof}
The source $A(n)$ is a dyadic unit for every nonnegative integer $n$.
Each nonempty summand has $m\ge M+1$, so
\eqref{eq:root:weights} and integrality of the binomial coefficient give
the asserted bound. At $N=M+1$, only the first omitted summand remains.
Its binomial coefficient and $A(0)$ both equal one.
\end{proof}

Thus later row-dependent cancellation can improve a particular tail,
but cannot increase the uniform exponent required to hold at every row.

\subsection{A finite carry minimum and an exact all-row factorization}
\label{app:root:factorization}

Let $\mathbb Z_2\langle x\rangle$ denote the restricted power series:
their integral coefficients tend to zero dyadically. Extend $A$ by
\[
 \mathscr A(x)=\frac{9^{x+1}-1}{8(x+1)(2x+1)},
 \qquad \mathscr A(-1)=-\frac{\log9}{8}.
\]
The apparent singularity at $x=-1$ is removable. Since
$v_2(\log9)=3$, its convergent expansion is
\begin{equation}
 \mathscr A(x)=\frac{\log9/8}{1+2x}
   \sum_{h\ge0}\frac{(\log9)^h}{(h+1)!}(x+1)^h.
 \label{eq:root:source}
\end{equation}
It follows coefficientwise that
$\mathscr A\in\mathbb Z_2\langle x\rangle^\times$,
$\mathscr A\equiv1\pmod2$, and $\mathscr A(N)=A(N)$.

Set $r=M+1\ge1$ and define
\[
 P_q(x)=\prod_{i=0}^{q-1}(x-i)(2x-(2i+1)),\qquad P_0=1,
\]
\[
 H_r(x)=\sum_{q\ge0}8^q\frac{\beta_{r+q}}{\beta_r}
                    P_q(x)\mathscr A(x-q).
\]
The coefficient valuations tend to infinity. At any nonnegative integer
$x$, terms with $q>x$ vanish. Cancelling the factorials in
\eqref{eq:root:packet} therefore gives the exact identity
\begin{equation}
 \mathcal T_{r-1}(N)=\kappa_r\binom{2N}{2r}H_r(N-r)
 \qquad(N\ge r).\label{eq:root:factor}
\end{equation}
This is a finite rational identity at every original Euler row.

The scalar packet weights have exact valuations
\[
 \lambda_r(q)=3q+s_2(r+q)-s_2(r).
\]
Put
\[
 \mu_r=\min_{q\ge0}\lambda_r(q),\qquad
 J_r=\{q:\lambda_r(q)=\mu_r\},\qquad G_r=2^{-\mu_r}H_r.
\]
Since $\lambda_r(0)=0$ and
$\lambda_r(q)\ge3q+1-s_2(r)$, every minimizing address satisfies
\begin{equation}
 0\le q\le\left\lfloor\frac{s_2(r)-1}{3}\right\rfloor.
 \label{eq:root:cutoff}
\end{equation}
In particular $\mu_r\le0$, and the minimum is found by a finite carry
calculation. Reduction of the whole series gives
\begin{equation}
 \overline G_r(X)=\sum_{q\in J_r}X^{\underline q}
       \quad\text{in }\mathbb F_2[X].
 \label{eq:root:reduction}
\end{equation}
Its highest degree is $\max J_r$, with leading coefficient one.
Thus the normalization is the actual Gauss normalization.

\begin{theorem}[Exact classical tail valuation]
For every $r\ge1$ and $N\ge r$,
\begin{equation}
 v_2\bigl(\mathcal T_{r-1}(N)\bigr)
 =4r-1+v_2\binom Nr+\mu_r+v_2G_r(N-r).
 \label{eq:root:exactvaluation}
\end{equation}
There are unique factors
\[
 G_r(X)=W_r(X)U_r(X),
\]
where $W_r\in\mathbb Z_2[X]$ is monic of degree $d_r=\max J_r$ and
$U_r\in\mathbb Z_2\langle X\rangle^\times$ satisfies $U_r\equiv1\pmod2$.
Consequently $v_2G_r(x)=v_2W_r(x)$ for every $x\in\mathbb Z_2$, and
\[
 \deg W_r\le\left\lfloor\frac{s_2(r)-1}{3}\right\rfloor.
\]
\label{thm:root:finitepreparation}
\end{theorem}
\begin{proof}
Equation~\eqref{eq:root:factor}, the exact weight in
\eqref{eq:root:weights}, and the digit-sum identity
$v_2\binom{2N}{2r}=v_2\binom Nr$ prove
\eqref{eq:root:exactvaluation}. The binomial coefficients themselves
need not be equal.

For preparation, start from the monic polynomial
$g=\overline G_r$ and its zero-one lift $W_1$, with $U_1=1$.
Suppose $G_r-W_kU_k$ is coefficientwise divisible by $2^k$.
Its divided reduction is a finite polynomial because $G_r$ is restricted.
Divide it by $g$ in $\mathbb F_2[X]$:
\[
 \overline{(G_r-W_kU_k)/2^k}=gq+s,\qquad\deg s<d_r.
\]
With zero-one coefficient lifts, set
$W_{k+1}=W_k+2^ks$ and $U_{k+1}=U_k+2^kq$.
The error is now divisible by $2^{k+1}$, the polynomial remains monic of
degree $d_r$, and the unit reduces to one. The sequences converge in the
Gauss norm. Uniqueness of quotient and remainder at successive binary
digits proves uniqueness of the factors. For $d_r=0$, take $W_r=1$.
The degree bound follows from \eqref{eq:root:cutoff}.
\end{proof}

This theorem reduces every depth to a finite polynomial without discarding
the infinite analytic tail. It does not explicitly classify every
$W_r$, decide every repeated root from a fixed precision, or prove that
all its dyadic roots are simple.

\subsection{Regular walls: exact root emergence}
\label{app:root:regular}

For a wall label $j\ge1$ and $0\le q<j$, define
\[
 e_j(q)=3q+s_2(j-1)-s_2(j-q-1),\qquad
 E_j=\min_{q<j}e_j(q),
\]
\[
 \Theta(j)=3j+1+s_2(j-1),\qquad L_j=\Theta(j)-E_j.
\]
Here $E_j$ is a carry minimum, not an Euler number. A wall label is
\emph{regular} when $e_j(q)>0$ for $1\le q<j$.
For $r=2^L-j$ with $L>L_j$, put $\delta=L-L_j$.
The carry identities give
\[
 \mu_r=E_j-\delta,\qquad J_r=\{j\},\qquad
 G_r(x)=P_j(x)U(x-j)+2^\delta V(x),
\]
where $U\equiv1\pmod2$. For a regular label, also
$E_j=0$ and $V\equiv1\pmod2$.

For $0\le i<j$, let
\[
 d_i=v_2(i!)+v_2((j-1-i)!),\qquad
 k_i=1+\left\lfloor\log_2\max(i,j-1-i)\right\rfloor,
\]
\[
 B_i=i+2^{k_i}\mathbb Z_2,\qquad t_i=d_i+k_i.
\]
At $j=1$ use $d_0=k_0=t_0=0$. The balls $B_i$ partition
$\mathbb Z_2$. Indeed a residue represented twice among $0,\ldots,j-1$
splits into its two children. Every other represented residue remains
one ball. On each such ball,
\[
 P_j(i+2^{k_i}y)=2^{t_i}yQ_i(y),\qquad
 Q_i\in\mathbb Z_2\langle y\rangle^\times,
 \qquad Q_i\equiv1\pmod2.
\]

\begin{theorem}[Regular-wall root law]
For every regular label $j$ and every integer $\delta\ge1$,
\[
 v_2G_r(x)=
 \begin{cases}
  \delta,&x\in B_i,\ \delta<t_i,\\
  d_i+v_2(x-\alpha_i),&x\in B_i,\ \delta\ge t_i.
 \end{cases}
\]
In the second case there is exactly one root $\alpha_i$ in $B_i$,
it is simple, and $v_2(\alpha_i-i)=\delta-d_i$.
Thus the number of dyadic roots is exactly
\[
 \#\{i:0\le i<j,\ t_i\le\delta\}.
\]
\label{thm:root:regular}
\end{theorem}
\begin{proof}
If $\delta<t_i$, the unit perturbation has smaller valuation than
the skeleton term throughout $B_i$, so the valuation is exactly
$\delta$. Otherwise, after the ball substitution and division by
$2^{t_i}$, the coefficientwise reduction is $y+c$ for a constant
$c\in\mathbb F_2$. A restricted integral series with this reduction
is a dyadic isometry: division of its difference by $y-z$ gives an
odd linear coefficient plus even higher contributions. Binary lifting
therefore gives its unique root and exact distance law. Evaluating at
$x=i$ gives the stated displacement, since
$G_r(i)=2^\delta V(i)$ has exact valuation $\delta$.
The partition proves exhaustion of the roots.
\end{proof}

The largest emergence threshold is
\[
 \max_i t_i=v_2((j-1)!)+\lceil\log_2 j\rceil.
\]
For $j=2^a$, all thresholds equal $j-1$. There are no roots before
that distance and exactly $j$ simple roots at and after it, within the
positive-distance range. The equality and pre-wall cases are also
resolved: regular labels $j\ge2$ have no roots there. $j=1$ has one
simple odd root at equality.

For example, $j=5$ has $(t_0,\ldots,t_4)=(6,3,4,3,6)$.
The exact root counts at distances $1,2,3,4,5,\ge6$ are
$0,0,2,3,3,5$. The older sufficient stages $3,5,7$ are not the
sharp emergence stages, which are $3,4,6$.
The first irregular label, $j=17$, has $V\equiv1+X\pmod2$.
The regular-wall theorem does not apply. Complete coefficient
certificates close distances $1$ through $19$, and the general separation
theorem gives all seventeen simple roots for $\delta\ge20$.
Those two kinds of evidence together resolve that family.

\subsection{What the later simplicity reductions prove}
\label{app:root:families}

For a general depth $r$, call $j_r=\max J_r$ its \emph{true terminal}.
This definition depends on the actual carry minimum. Rewriting the same
integer as another wall expression does not change $j_r$.
The universal simplicity target is
\begin{equation}
 \text{for every integer }r\ge1\text{ and }x\in\mathbb Z_2,
 \qquad G_r(x)=0\ \Longrightarrow\ G_r'(x)\ne0.
 \label{eq:root:simplicitytarget}
\end{equation}
The following are completed restricted results of the later continuations.
Their family statements include a proof that the indicated label is the
true terminal throughout the displayed ranges.

\begin{center}
\begin{tabular}{lll}
\toprule
Depth $r$ & Odd multiplier range & Outer scales\\
\midrule
$B2^L-129$ & $0<B<2^{89}$ & $L\ge392$\\
$B2^L-161$ & $0<B<2^{932}$ & $L\ge487$\\
\bottomrule
\end{tabular}
\end{center}
All dyadic roots are simple in both infinite integer families.
For arbitrary larger admissible multipliers, a multiple root with true
terminal $129$ can occur only at $L=515$ and in two necessary multiplier
classes modulo $2^{90}$. With true terminal $161$, it can occur only at
$L=640$ and in one necessary multiplier class modulo $2^{935}$.
These classes are necessary conditions. The records do not assert
existence of a singular ordinary Euler index in them.

For terminal $161$, the neighboring phases are resolved uniformly:
$L=639$ has exactly one simple root and $L=641$ exactly three simple
roots, for every positive odd multiplier with that true terminal.
The middle phase is reduced, before the deeper certificate is applied,
to
\[
 B\equiv17659\pmod{2^{18}},\qquad
 x\equiv10202\pmod{2^{15}}.
\]
The $935$-bit multiplier condition composes this whole-phase reduction
with an inherited $917$-bit local parameter certificate. It is neither
a new $935$-bit reconstruction nor an all-precision interpolation of
the undepleted Euler family.

The structural argument also extends beyond lower address one.
For a certified true-terminal wall, let $a$ be its last lower minimizing
address, $k=j_r-a$ the regular complement, and $t=v_2(k)$.
Under the regular lower-address hypotheses of the source theorem,
all roots above emerged complement balls are simple. Every root is
simple at every wall distance in the band
\[
 \Theta(a)\le t<\Theta(a)+v_2(a!).
\]
Equality walls are simple even outside that band. This includes
certified terminal-$258$ walls with $a=2$, $k=256$, and $L\ge777$.
For $a=3$, $k=4096$, only $L=16385$ remains a possible critical scale.
Outside the band, the theorem gives necessary local valuation conditions.
It does not exclude the remaining unit cancellation. Irregular lower
addresses need their own prefix analysis.

\subsection{The actual stationary discriminant}
\label{app:root:stationary}

The current focused obstruction is the admitted ordinary-index family
\begin{equation}
 r=(17659+2^{18}T)2^{640}-161,
 \qquad T\in\mathbb Z_{\ge0},\qquad j_r=161.
 \label{eq:root:focusedfamily}
\end{equation}
On its surviving chart define the \emph{actual}, undepleted function
\[
 K_r(Z)=2^{-178}G_r(10202+2^{15}Z).
\]
The complete coefficient theorem gives
\begin{equation}
 K_r(Z)=a_r+2b_rZ+u_rZ^2+256Z^3R_r(Z),
 \quad a_r,b_r\in\mathbb Z_2,\quad
 u_r\in\mathbb Z_2^\times,\quad
 R_r\in\mathbb Z_2\langle Z\rangle.
 \label{eq:root:quadraticclass}
\end{equation}
Its half derivative is a dyadic isometry. It therefore has a unique
critical point $c_r\in\mathbb Z_2$. Set
$\Phi_r^{\mathrm{act}}=K_r(c_r)$.
The earlier independent-undepletion notation for this same scalar is
$\Phi_r(2^{2r-25})$.
Excluding a multiple root on the surviving chart is exactly the assertion
$\Phi_r^{\mathrm{act}}\ne0$.

\begin{proposition}[Exact quadratic compression]
There is a unique preparation
\[
 K_r(Z)=\mathcal U_r(Z)\mathcal W_r(Z),\qquad
 \mathcal W_r(Z)=Z^2+2\beta_rZ+\alpha_r,\qquad
 \mathcal U_r\equiv u_r\pmod{256}.
\]
With $\delta_r=\beta_r^2-\alpha_r$ and
$\tau_r=\mathcal U_r'(c_r)/\mathcal U_r(c_r)$,
\begin{equation}
 \delta_r=-\frac{\Phi_r^{\mathrm{act}}}{\mathcal U_r(c_r)}
  \left(1-\frac{\tau_r^2\Phi_r^{\mathrm{act}}}{4\mathcal U_r(c_r)}\right),
 \qquad v_2(\delta_r)=v_2(\Phi_r^{\mathrm{act}}).
 \label{eq:root:discriminant}
\end{equation}
The equality includes the zero case.
\end{proposition}
\begin{proof}
The same successive monic polynomial division used in
Theorem~\ref{thm:root:finitepreparation}, started at precision eight,
preserves the quadratic and unit reductions of
\eqref{eq:root:quadraticclass}. Write
$W=Z^2+2\beta_rZ+\alpha_r$ and $w=W(c_r)$.
Stationarity gives $W'(c_r)=-\tau_rw$, while
$w=\Phi_r^{\mathrm{act}}/\mathcal U_r(c_r)$.
The quadratic identity
$\beta_r^2-\alpha_r=W'(c_r)^2/4-W(c_r)$ proves the formula.
Since $\tau_r\in256\mathbb Z_2$, its bracket is a unit.
\end{proof}

This compression retains the unit-cofactor derivative. It does not
assert that $c_r$, $\alpha_r$, or $\beta_r$ is rational or algebraic
over $\mathbb Q$. The remaining arithmetic obligation is therefore
\begin{equation}
 \delta_r\ne0\quad\text{for every admitted ordinary index in
 \eqref{eq:root:focusedfamily}.}
 \label{eq:root:remaining}
\end{equation}
Even this assertion would still leave the terminal-$129$ and higher-address
obligations in the universal target \eqref{eq:root:simplicitytarget}.

\subsection{Approximation reaches every depth, but separation remains open}
\label{app:root:approximation}

The latest continuation retains the exact Bernoulli undepletion factor
\[
 e_{r,q}=\frac{1-2^{2r-1}}{1-2^{2r-1}4^q},\qquad q\ge0.
\]
It is a dyadic unit, including $e_{r,0}=1$. Multiplication by this
factor preserves the all-$q$ error of the shifted Jacobi moment model.
Consequently the actual model critical values satisfy, for every $m\ge25$,
\begin{equation}
 v_2\bigl(\Phi_{r,m}^{\mathrm{act}}-\Phi_r^{\mathrm{act}}\bigr)
 \ge L_m=8m-4s_2(m)-178,
 \qquad L_m\longrightarrow\infty.
 \label{eq:root:actualapprox}
\end{equation}
The estimate includes every packet and controls motion of the critical
point. There is no ceiling at the undepletion exponent $2r-25$.
For each fixed admitted $r$ it implies the exact equivalence
\begin{equation}
 \Phi_r^{\mathrm{act}}\ne0
 \quad\Longleftrightarrow\quad
 \text{there exists }m\ge25\text{ such that }
 v_2(\Phi_{r,m}^{\mathrm{act}})<L_m.
 \label{eq:root:stricttest}
\end{equation}
Indeed a passing strict test has the same valuation as the limit. Conversely
any finite limit valuation is eventually smaller than $L_m$.
The existence of a passing depth for every ordinary index is unproved.
Nonzero models at every depth are compatible with an exactly zero limit.

The degree-two preparation also reduces the old representative-height
cost to a linear cost, about $2P$ bits at precision $P$. This remains
larger than the $P$ bits needed for a strict nonzero residue certificate.
The fixed rational control $(Z-1/3)^2$ has zero critical value while
its least-residue integer discriminant representatives are nonzero at
every even $P\ge6$. Thus a nonzero integer representative alone cannot
prove a nonzero dyadic discriminant.

One sufficient new arithmetic input would be an odd-denominator
projective quadratic
\[
 A Z^2+2BZ+C\equiv A\mathcal W_r(Z)\pmod{2^P},\qquad A\text{ odd},
 \qquad 0<|B^2-AC|<2^P,
\]
at a proved precision for every admitted ordinary $r$.
Then $(B^2-AC)/A^2\equiv\delta_r\pmod{2^P}$ cannot vanish
modulo $2^P$. No depth-independent theorem supplying this inequality
has been established. The tested recurrence and exterior-contact
identities retain balancing terms and do not force a contradiction.
Finite exponential models approximate the actual target at every depth,
but their order grows with precision, their clearing polynomial has
zeros on the actual chart, and their proved construction-specific height
cost does not give the required discriminant bound.

\subsection{Status and source map}
\label{app:root:evidence}

\begin{center}
\begin{tabular}{>{\raggedright\arraybackslash}p{0.64\linewidth}>{\raggedright\arraybackslash}p{0.26\linewidth}}
\toprule
Claim & Status at this edition\\
\midrule
Nishibuchi Conjecture~3.9 in its stated parameter range
 & Resolved in the main text\\
Sharp uniform ordinary Euler tail exponent $4M+3$
 & Proved, with earlier formal coverage\\
Exact tail factorization and finite polynomial at every depth
 & Proved in the later prose development\\
Regular-wall emergence and indicated integer simplicity families
 & Proved within their stated hypotheses\\
Actual quadratic compression and unbounded stationary approximation
 & Proved. No fresh Lean build\\
Nonzero actual stationary discriminant for every admitted index
 & Open\\
Universal simplicity of every $G_r$
 & Open\\
Native~13.2 presentation invariance and Native~13.3 global/completed gluing
 & Separate open general programs\\
\bottomrule
\end{tabular}
\end{center}

The tail factorization, regular-wall proofs, finite preparation, and
coefficient cutoffs are recorded in the preserved
\nolinkurl{sources/euler_v4.md} of the regular-address checkpoint.
The two integer families and multiplier conditions are proved in its
\nolinkurl{next/terminal129/PROOF.md} and
\nolinkurl{next/critical_phases/PROOF.md}.
\nolinkurl{local_parameter/RECENTER_PROOF.md} supplies the inherited
local precision certificate. The higher-address extension is in
\nolinkurl{next/higher_address/regular_address_transversality.md}.
The stationary checkpoint records the exact discriminant identity in
\nolinkurl{stationary_separation/ordinary/QUADRATIC_COMPRESSION_AND_HEIGHT.md}
and actual convergence in
\nolinkurl{stationary_separation/root/ACTUAL_LIMIT_AND_SEPARATION.md}.
Its synthesis and proof reviews retain the failed
implications and the precise remaining arithmetic input.

The source archive's \nolinkurl{evidence/root_scope_inventory.md} gives
the full checkpoint paths and content hashes used for this appendix.
The 2026-10-05 complete-work archive preserves the original layouts and
historical releases. Its stored hashes and verification receipts are
evidence of integrity and recorded execution. They are not substitutes
for a proof of the open assertions above.
